\section{Emission source identification}

The peak detector scan showed several frequencies close to or above the limit line.
However, a peak detector result does not by itself prove that the final quasi peak result
will exceed the limit. This can be seen in the final measurement results, where only one
of the selected frequencies exceeded the limit after quasi peak measurement. Performing
a full quasi peak measurement at every peak in the spectrum was not feasible within the
scope of this study. Therefore, the source identification is focused first on the frequency
that was confirmed to exceed the limit, namely \(480\,MHz\).

\subsection{480 MHz}

The \(480\,MHz\) emission is interesting because it is the second harmonic of \(240\,MHz\),

\[
480\,\mathrm{MHz} = 2 \cdot 240\,\mathrm{MHz}.
\]

This is relevant because \(240\,MHz\) is the CPU clock frequency used by the ESP32. The
spectrum also shows peaks close to \(719\,MHz\) and \(970\,MHz\), which are close to the
third and fourth harmonics of \(240\,MHz\),

\[
3 \cdot 240\,\mathrm{MHz} = 720\,\mathrm{MHz},
\]

\[
4 \cdot 240\,\mathrm{MHz} = 960\,\mathrm{MHz}.
\]

There is also a visible peak around \(240\,MHz\) in the average trace, although it is not high
enough to exceed the limit. This suggests that the \(240\,MHz\) clock family may be present
in the measured spectrum, but it does not necessarily mean that the ESP32 clock itself is
the complete problem.

If the ESP32 clock alone was enough to create this level of radiated emission, the same
problem would likely be common in many products using the same module. It is therefore
more reasonable to suspect that something in the AGV design allows this clock-related
energy to couple into a structure that radiates efficiently. This could be a cable, a supply
path, a ground return path, or another conductive part of the robot platform.

One possible explanation is that one or more cables act as unintended antennas. As the
frequency increases, the wavelength decreases according to

\[
\lambda = \frac{c}{f}.
\]

For \(480\,MHz\) this gives

\[
\lambda = \frac{3 \cdot 10^8}{480 \cdot 10^6}
          \approx 0.625\,\mathrm{m}.
\]

A quarter wavelength is therefore

\[
\lambda/4 \approx 0.156\,\mathrm{m}.
\]

This is approximately \(16\,cm\), which is comparable to cable lengths and wiring sections in
the AGV. Therefore, a cable does not need to be very long to become an efficient
radiating structure at \(480\,MHz\). The \(480\,MHz\) peak can therefore be interpreted as a
possible clock-related disturbance that becomes critical because of the physical AGV
implementation, rather than as proof that the ESP32 module alone is the source.
